\textbf{Phần 1. Trắc nghiệm lựa chọn}
\setcounterpageref{ex}{0}

\begin{ex} %[2D4V23]
Tính kết quả của phép trừ: $72 - 38$.
\choice
{$33$}
{$35$}
{\True $34$}
{$44$}
\loigiai{
$72 - 38 = 34$.
\textbf{Đáp án của câu hỏi là C}
}
\end{ex}

\begin{ex} %[2D4V23]
Một mảnh vườn có $85$ cây bưởi và cây cam, trong đó có $48$ cây bưởi. Hỏi có bao nhiêu cây cam?
\choice
{$36$ cây cam}
{\True $37$ cây cam}
{$38$ cây cam}
{$47$ cây cam}
\loigiai{
$85 - 48 = 37$ (cây cam).
\textbf{Đáp án của câu hỏi là B}
}
\end{ex}

\begin{ex} %[2D4V23]
Tìm $x$ biết: $x + 47 = 81$.
\choice
{$32$}
{$33$}
{$35$}
{\True $34$}
\loigiai{
$x = 81 - 47 = 34$.
\textbf{Đáp án của câu hỏi là D}
}
\end{ex}

\begin{ex} %[2D4V23]
Một cửa hàng có $90\text{ kg}$ gạo, buổi sáng bán được $36\text{ kg}$ gạo, buổi chiều bán được $28\text{ kg}$ gạo. Hỏi cửa hàng còn lại bao nhiêu ki-lô-gam gạo?
\choice
{\True $26\text{ kg}$}
{$24\text{ kg}$}
{$25\text{ kg}$}
{$28\text{ kg}$}
\loigiai{
Tổng gạo đã bán: $36 + 28 = 64\text{ (kg)}$. Còn lại: $90 - 64 = 26\text{ (kg)}$.
\textbf{Đáp án của câu hỏi là A}
}
\end{ex}

\begin{ex} %[2D4V23]
Tính giá trị của biểu thức: $64 - 27 - 18$.
\choice
{$18$}
{$20$}
{\True $19$}
{$17$}
\loigiai{
$64 - 27 = 37$; $37 - 18 = 19$.
\textbf{Đáp án của câu hỏi là C}
}
\end{ex}

\begin{ex} %[2D4V23]
Cho phép tính: $71 - \overline{a5} = 36$. Giá trị của chữ số $a$ là
\choice
{$2$}
{\True $3$}
{$4$}
{$1$}
\loigiai{
Ta có: $\overline{a5} = 71 - 36 = 35$. Do đó $a = 3$.
\textbf{Đáp án của câu hỏi là B}
}
\end{ex}

\begin{ex} %[2D4V23]
Một quãng đường dài $83\text{ km}$, ô tô đã đi được $49\text{ km}$. Hỏi ô tô còn phải đi bao nhiêu ki-lô-mét nữa để hết quãng đường?
\choice
{$32\text{ km}$}
{$33\text{ km}$}
{$35\text{ km}$}
{\True $34\text{ km}$}
\loigiai{
$83 - 49 = 34\text{ (km)}$.
\textbf{Đáp án của câu hỏi là D}
}
\end{ex}

\begin{ex} %[2D4V23]
Hiệu của số tròn chục lớn nhất có hai chữ số và số liền trước của $46$ là
\choice
{\True $45$}
{$44$}
{$46$}
{$43$}
\loigiai{
Số tròn chục lớn nhất có 2 chữ số là 90. Số liền trước 46 là 45.\nHiệu: $90 - 45 = 45$.
\textbf{Đáp án của câu hỏi là A}
}
\end{ex}

\begin{ex} %[2D4V23]
Điền dấu thích hợp: $82 - 37 \dots 71 - 25$.
\choice
{$>$}
{$=$}
{\True $<$}
{Không so sánh được}
\loigiai{
$82 - 37 = 45$; $71 - 25 = 46$. Vì $45 < 46$ nên điền dấu $<$. 
\textbf{Đáp án của câu hỏi là C}
}
\end{ex}

\begin{ex} %[2D4V23]
Hai số có tổng bằng $74$, biết số lớn là $48$. Số bé là
\choice
{$25$}
{\True $26$}
{$27$}
{$28$}
\loigiai{
Số bé: $74 - 48 = 26$.
\textbf{Đáp án của câu hỏi là B}
}
\end{ex}

\begin{ex} %[2D4V23]
Mẹ mua một tấm vải dài $60\text{ dm}$. Mẹ may váy hết $28\text{ dm}$, may áo hết $17\text{ dm}$. Tấm vải còn lại dài bao nhiêu đề-xi-mét?
\choice
{$14\text{ dm}$}
{$16\text{ dm}$}
{$13\text{ dm}$}
{\True $15\text{ dm}$}
\loigiai{
Mẹ may hết: $28 + 17 = 45\text{ (dm)}$. Còn lại: $60 - 45 = 15\text{ (dm)}$.
\textbf{Đáp án của câu hỏi là D}
}
\end{ex}

\begin{ex} %[2D4V23]
Tìm một số biết rằng lấy $84$ trừ đi số đó thì được số bé nhất có hai chữ số giống nhau.
\choice
{\True $73$}
{$72$}
{$74$}
{$71$}
\loigiai{
Số bé nhất có hai chữ số giống nhau là 11. Số đó là: $84 - 11 = 73$.
\textbf{Đáp án của câu hỏi là A}
}
\end{ex}

\textbf{Phần 2. Trắc nghiệm đúng sai}
\setcounterpageref{ex}{0}

\begin{ex} %[2D4V23]
Xét tính đúng sai của phép trừ có nhớ trong phạm vi 100:
\choiceTF
{\True Phép trừ $63 - 28$ có kết quả là $35$.}
{Phép trừ $70 - 35$ có kết quả là số chẵn.}
{\True Khi bớt ở cả số bị trừ và số trừ đi cùng một số đơn vị thì hiệu không đổi.}
{\True Hiệu của $81 - 49$ bằng hiệu của $82 - 50$.}
\loigiai{
\begin{itemize}
    \item Ý a) Đúng: $63 - 28 = 35$.
    \item Ý b) Sai: $70 - 35 = 35$, là số lẻ.
    \item Ý c) Đúng: Tính chất của phép trừ.
    \item Ý d) Đúng: Cùng thêm 1 đơn vị vào cả số bị trừ và số trừ: $81 - 49 = 82 - 50 = 32$. Đúng.
\end{itemize}
\begin{itemize}
    \item a) \textbf{Đúng}.
    \item b) \textbf{Sai}.
    \item c) \textbf{Đúng}.
    \item d) \textbf{Đúng}.
\end{itemize}
}
\end{ex}

\begin{ex} %[2D4V23]
Một kho thóc có 95 bao thóc. Ngày thứ nhất xuất 38 bao, ngày thứ hai xuất 29 bao:
\choiceTF
{\True Tổng số bao thóc đã xuất trong hai ngày là 67 bao.}
{\True Số thóc còn lại trong kho là 28 bao.}
{Ngày thứ nhất xuất ít hơn ngày thứ hai 9 bao.}
{Kho thóc còn lại nhiều hơn số thóc ngày thứ hai xuất.}
\loigiai{
\begin{itemize}
    \item Ý a) Đúng: $38 + 29 = 67$ bao.
    \item Ý b) Đúng: $95 - 67 = 28$ bao.
    \item Ý c) Sai: Ngày thứ nhất (38 bao) nhiều hơn ngày thứ hai (29 bao) là $38 - 29 = 9$ bao.
    \item Ý d) Sai: Còn lại 28 bao, ngày thứ hai xuất 29 bao nên còn lại ít hơn.
\end{itemize}
\begin{itemize}
    \item a) \textbf{Đúng}.
    \item b) \textbf{Đúng}.
    \item c) \textbf{Sai}.
    \item d) \textbf{Sai}.
\end{itemize}
}
\end{ex}

\textbf{Phần 3. Trả lời ngắn}
\setcounterpageref{ex}{0}

\begin{ex} %[2D4V23]
Tính kết quả: $83 - 47$.
\shortans{36}
\loigiai{
\textbf{Đáp số:} $36$\\
$83 - 47 = 36$.
}
\end{ex}

\begin{ex} %[2D4V23]
Một giá sách có 70 quyển sách, đã cho mượn 34 quyển. Hỏi trên giá sách còn lại bao nhiêu quyển sách?
\shortans{36}
\loigiai{
\textbf{Đáp số:} $36$\\
$70 - 34 = 36$ quyển.
}
\end{ex}

\begin{ex} %[2D4V23]
Tìm số tự nhiên y biết: $y + 38 = 91$.
\shortans{53}
\loigiai{
\textbf{Đáp số:} $53$\\
$y = 91 - 38 = 53$.
}
\end{ex}

\begin{ex} %[2D4V23]
Cuộn dây dài 92 cm, bạn Nam cắt đi một đoạn dài 56 cm. Hỏi đoạn dây còn lại dài bao nhiêu xăng-ti-mét?
\shortans{36}
\loigiai{
\textbf{Đáp số:} $36$\\
$92 - 56 = 36\text{ cm}$.
}
\end{ex}

\begin{ex} %[2D4V23]
Tìm hiệu của số lớn nhất có hai chữ số và số 67.
\shortans{32}
\loigiai{
\textbf{Đáp số:} $32$\\
$99 - 67 = 32$.
}
\end{ex}

\begin{ex} %[2D4V23]
Lớp 2A có 33 học sinh, lớp 2B có 30 học sinh. Cả hai lớp có 27 bạn nam. Hỏi cả hai lớp có bao nhiêu bạn nữ?
\shortans{36}
\loigiai{
\textbf{Đáp số:} $36$\\
Tổng số học sinh: $33 + 30 = 63$. Số bạn nữ: $63 - 27 = 36$ bạn.
}
\end{ex}
